\section{Data} \label{sec:data_appendix}


This section aims to discuss the various data sets employed within this work, giving for each one an overview of the dataset, a description of the source of the data, the methodological details of how the data is obtained and analyzed, possible systematic corrections to be applied, the most important assumptions going into the dataset or method, and the connection to the network. Note that this section primarily describes the data itself---for an overview over the respective working principles, see Section~\ref{sec:data_summary}, and for an overview of the implementation within the {\DN} set of equations, see Appendix~\ref{app:equations}.

\subsection{Geometric distances} \label{sec:anchors}

\subsubsection{Distances to Galactic Cepheids from trigonometric parallaxes}

\paragraph{Overview.} 
Trigonometric parallaxes of Galactic Cepheids were measured with high signal-to-noise by the ESA \textit{Gaia} mission \citep{2016A&A...595A...1G}. The \textit{Gaia} spacecraft was in orbit around the Lagrange point L2 and traveled together with the Earth around the Sun once per year. Given {\it Gaia}'s distance from the Earth and Earth's distance to the Sun, {\it Gaia}'s large baseline made it extremely sensitive and allowed it to measure, with observations of the same star collected over a few years, the apparent shift (parallax) of stars that were within $d < 10$ kpc.

\paragraph{Data Sources.}
Parallaxes of Cepheids measured by the ESA {\it Gaia} mission are published as part of the early third data release~\citep[GEDR3,][]{2021A&A...649A...1G,2021A&A...649A...2L}. Seven additional narrow-angle parallaxes of Milky Way Cepheids measured by spatially scanning {\it HST}~\citep{2014ApJ...785..161R,2016ApJ...825...11C} were added from~\citet{2018ApJ...855..136R}. 

\paragraph{Methodological Details.}
``Field'' Cepheids identify classical Cepheids in the Milky Way with direct trigonometric parallax measurements. Cepheids in clusters have also been the focus of multiple studies: since {\it Gaia} provides parallaxes for all individual members, averaging them yields a value nearly three times more precise than for field Cepheids~\citep{2020A&A...643A.115B, 2022ApJ...938...36R, 2023A&A...672A..85C, 2024AJ....168...34W}. 

\paragraph{Systematic Corrections.}
Well documented systematics of GEDR3 parallaxes were corrected using the formalism by \citet[L21]{2021A&A...649A...4L}, which considers source magnitude, color, and on-sky position (sine of ecliptic latitude). It is also well documented that additional corrections (known as a residual parallax offset $\varpi_{\rm corr}$) are required at a sample level to ensure unbiased parallaxes~\citep[e.g.,][]{2023A&A...680A.105K,2022AJ....163..149W,2021ApJ...909..200B}.~\citet{2021ApJ...908L...6R} simultaneously solved for the fiducial absolute magnitude of Cepheids and for $\varpi_{\rm corr}$ for a sample of 75 field Cepheids observed photometrically with {\it HST} in the SH0ES Wesenheit index (a color-corrected magnitude, $m_H^W$) and found $\varpi_{\rm corr}=-14\,\mu$arcsec. For cluster members (used to infer parallaxes of cluster Cepheids),~\citet[R22]{2022ApJ...934L...7R} find no evidence for a residual parallax offset: the corrections from~\citet{2021A&A...649A...4L} effectively mitigate the parallax bias, as cluster members fall within the color and magnitude range of the sources used in its calibration.  The cluster parallax uncertainty is a combination of the statistical uncertainty and the small-scale angular covariance, with the latter being the dominant term. Table 1 of~\citet{2022ApJ...938...36R} presents the mean parallaxes of 17 Cluster Cepheids from GEDR3 also observed photometrically with HST.

\paragraph{Assumptions.}
For field Cepheids, we assume that the parallax as measured by GEDR3 has a normal distribution with central value equal to the reciprocal of the true distance plus the fixed offset $\varpi_{\rm corr} $ discussed above, and width given by the nominal EDR3 parallax uncertainty for each source.  The calibration of Milky Way Cepheids is obtained from the aggregation of all the measured parallax values, including parallax, photometric, and systematic uncertainties, and marginalized over the solution of the offset value.  The uncertainty in the offset $\varpi_{\rm corr} $ thus determined is therefore included in the final calibration uncertainty.  
For cluster Cepheids, the distance to each cluster is estimated from the aggregation of the parallax measurements for cluster stars; as noted above, cluster stars are generally within the applicability parameters of the standard offset determined in~\citet{2021A&A...649A...4L}.  However, it is important to note that parallax uncertainties within a small angular region are significantly correlated; this correlation is included in the estimate of the combined uncertainty in~\citet{2022ApJ...938...36R}.  Thanks to the combination of multiple sources, individual cluster parallaxes are typically a factor of three better than the nominal uncertainty for a single bright Cepheid.

\paragraph{Connection to the Network.}
Trigonometric parallaxes from {\it Gaia} GEDR3 are used as an anchor of the distance ladder and calibrate the Cepheid period-luminosity relation within the Milky Way.


\subsubsection{Distance to NGC 4258 using Megamasers}
\label{sec:NGC4258MM}

\paragraph{Overview.}
NGC 4258 is a nearby Seyfert 2 galaxy that is well-known for hosting the archetypal circumnuclear water megamaser disk system~\citep{1988IAUS..129..231C,1995ApJ...440..619G}.  VLBI mapping and spectral monitoring of the water masers in this system reveal that they reside in an edge-on disk and appear to be executing nearly perfectly Keplerian orbits within the centermost $\sim$1\,pc of the galaxy, where the gravitational potential is dominated by the central SMBH~\citep{1995Natur.373..127M,1999Natur.400..539H}.  By modeling these maser orbits across $\sim$3.5 years of VLBI monitoring data~\citep[collected by][]{2007ApJ...659.1040A,2008ApJ...672..800H,2013ApJ...775...13H},~\citet{2019ApJ...886L..27R} determined a distance to NGC 4258 of $7.576$\,Mpc, with a statistical uncertainty of $0.082$\,Mpc and a systematic uncertainty of $0.076$\,Mpc.  For the analyses presented in this paper, we adopt a distance modulus of $\mu_{\text{N4258}} = 29.397 \pm 0.032$ for this system \citep{2019ApJ...886L..27R}.

\paragraph{Data Sources.}
The NGC 4258 data used by~\citet{2019ApJ...886L..27R} come from 18 epochs of VLBI monitoring with the Very Long Baseline Array, spanning 3.5 years and described by~\citet{2008ApJ...672..800H} and~\citet{2013ApJ...775...13H}.

\paragraph{Methodological Details.}
The orbital motion of the masers is determined using a combination of spectral and spatial information from the VLBI monitoring, which yields constraints on the sky-plane components of the masers' positions as well as the line-of-sight components of their velocities and accelerations.  The orbital parameters for all masers are modeled alongside global parameters describing the disk morphology, the mass and location of the central black hole, its line-of-sight velocity, and the distance to the system.

\paragraph{Systematic Corrections.}
None.

\paragraph{Assumptions.}
It is assumed that the masers execute Keplerian orbits within a thin disk and that their orbital motion is entirely determined by the gravity of an enclosed point mass. Additionally, the analysis assumes no significant evolution of fundamental physical constants (e.g., $G$, $c$, $\alpha$) across space, time, or environment.

\paragraph{Connection to the Network.}
The distance to NGC 4258 is used as an anchor in the network, where it calibrates the Cepheid variable, TRGB, JAGB, and Mira variable methods.


\subsubsection{Distances to the Magellanic Clouds from detached eclipsing binaries}
\label{sec:MagClDEB}

\paragraph{Overview.}
The eclipsing binary (EB) method is a spectrophotometric method that provides near-geometrical distances to detached EB stars. The method serves both as a standard ruler, via absolute radii of EB components, and as a standard candle, by means of the empirical surface brightness -- color relations (SBCR). The method is applied to two closest galaxies: 
the Large Magellanic Cloud~\citep{2019Natur.567..200P}, giving a distance of 49.59 kpc, and the Small Magellanic Cloud~\citep{2020ApJ...904...13G}, giving a distance of 62.43 kpc. Both galaxies serve as geometrical anchors of the extragalactic distance ladder. For the analysis in the present work we adopt distance moduli of $18.477\pm 0.004$ (stat.)$ \pm 0.026$ (sys.) mag and $18.977 \pm 0.016 $ (stat.)$ \pm 0.028$ (sys.) mag for the LMC and the SMC, respectively.

\paragraph{Data Sources.}
Calibration of SBCR was based on near-infrared interferometry measurements of angular diameters of 40 nearby red clump stars~\citep{2019A&A...632A..31G}. Their optical photometry came from compilation by~\citep{1997A&AS..124..349M} and near-infrared magnitudes came after~\citep{2012MNRAS.419.1637L}. Twenty detached EBs containing non-active late-type giant stars were analysed in the LMC~\citep{2018ApJ...860....1G} and 15 similar systems in the SMC~\citep{2020ApJ...904...13G}. Data used in both papers comprised 20 years of ground-based optical photometry from the OGLE project~\citep{1997AcA....47..319U}, ground-based NIR photometry secured with SOFI instrument on NTT telescope in La Silla Observatory and long-term spectroscopic follow-up with 6 and 8 meter class telescopes in Las Campanas Observatory and Paranal Observatory carried out by the Araucaria project~\citep[e.g.][]{2009ApJ...697..862P}.

\paragraph{Methodological Details.}
Light curves and radial velocity curves were analysed simultaneously with the updated Wilson-Devinney code~\citep{1971ApJ...166..605W} in order to derive masses and radii of eclipsing binary components with a precision better than 3 
percent. Interstellar extinction corrected  photometric indices were used to derive angular diameters of components utilizing the SBCR calibrated by~\citep{2019Natur.567..200P}. The geometric distance to a particular EB results from dividing a radius of a component by its angular diameter. Reddenings were determined by three independent methods using reddening maps~\citep{1998ApJ...500..525S,2011AJ....141..158H}, the sodium doublet~\citep{1997A&A...318..269M} and utilizing extinction-free photometric indices derived from analysis of decomposed spectra of components. 

\paragraph{Systematic Corrections.}
None.

\paragraph{Assumptions.}
The binary stars are assumed to follow Roche-lobe geometry in the modeling. No significant dependence of the SBCR on metallicity is assumed for G- and K-type giant stars. A standard interstellar extinction curve is used~\citep{2007ApJ...663..320F} with $R_V = 3.1$. Grey extinction is assumed to be negligible.

\paragraph{Connection to the Network.}
Distances to the Magellanic Clouds are used to calibrate the period-luminosity relation (Leavitt Law) for the classical Cepheids within these galaxies. 

\subsection{Stellar standard candles}

\subsubsection{Cepheid Variables}\label{ssec:cepheid}

\paragraph{Overview.}
Cepheid type variables are yellow supergiants in the instability strip that follow a period-luminosity relation, also called the Leavitt Law~\citep{1912HarCi.173....1L}. They are considered primary distance indicators and are commonly used to calibrate Type Ia supernovae (though in this paper their use extends to other secondary distance indicators).

\paragraph{Data Sources.}
Distances measured from Cepheids are adopted from the following sources:
\begin{itemize}
\item HST---Cepheids---SH0ES: All Cepheids in the four anchor galaxies were observed in the same three filters ($F160W$, $F555W$, $F814W$) with the Wide Field Camera 3 (WFC3) on the Hubble Space Telescope (HST). A first study including 7 Milky Way field Cepheids with spatial scanning parallaxes and photometry were presented by~\citet{2018ApJ...855..136R}. Later, the photometry for a sample of 75 Milky Way Cepheids with accurate GEDR3 parallaxes was obtained using the same spatial scanning mode to avoid saturation and is listed in Table 1 from~\citet{2021ApJ...908L...6R}.  The photometry for an additional 17 Milky Way Cepheids located in open clusters can be found in Table 2 from~\citet{2022ApJ...938...36R}. In the LMC, 70 Cepheids were measured in Table 2 from~\citet{2019ApJ...876...85R} in the same three filters. The SMC has a non-negligible depth along the line of sight, therefore only Cepheids in the core regions were observed: a sample of 88 Cepheids with HST photometry is given in Table 2 from~\citet{2024ApJ...973...30B}. Finally, in the maser host galaxy NGC 4258, HST photometry of 669 Cepheids in the optical ($F555W$, $F814W$) comes from Table 2 from~\citet{2022ApJ...940...64Y} and in the infrared ($F160W$) from~\citet{2022ApJ...934L...7R}. HST photometry for Cepheids in 37 SN~Ia host galaxies on the second rung are from~\citet{2022ApJ...934L...7R}.
\item JWST---Cepheids---SH0ES: Due to the brightness of Cepheids and the high sensitivity of JWST, the most practical and valuable set of measurements involves following up known Cepheids (discovered in the optical by HST) with JWST NIRCAM in the hosts of SNe~Ia and the maser host NGC 4258. Table A2 of \citet{2024ApJ...977..120R} presents distances to 8 SN Ia hosts measured following this approach.  See~\citet{2023ApJ...956L..18R} for details of how these observations were made.
\end{itemize}

\paragraph{Methodological Details.}
Cepheid pulsation periods are measured from HST optical broad band light curves (in host galaxies) and from well covered ground based light curves in the Milky Way and Magellanic Clouds.  The brightness of each Cepheid is measured from optical and IR observations using HST or JWST, and epoch-corrected using the estimated amplitude and phase of the light curve obtained from optical data.  Periods are known at the $10^{-4}$ day level for Cepheids in the MW and Magellanic Clouds and to 1--2\% precision for those determined with HST in more distant hosts. A linear relation is fitted between the logarithm of the pulsation period and Cepheid brightness, which gives the Period-Luminosity slope and intercept. Intercept differences between the period-luminosity relations in two galaxies give their relative distance.

\paragraph{Systematic Corrections.}
The intercept of the period-luminosity relation depends on the metallicity of Cepheids. Metallicity effects were incorporated by differences in oxygen abundances. Milky Way, LMC, and SMC oxygen abundances of Cepheids were measured using high-resolution spectroscopy~\citep{2022A&A...658A..29R, 2023ApJ...955L..13B, 2024ApJ...973...30B}. Oxygen abundances of Cepheids in host galaxies were estimated using oxygen abundance gradients determined via the $R_{23}$ strong-line method as described by \citet[][Sect.\,3.5]{2022ApJ...934L...7R} and the galactocentric distances of each Cepheid. The period-luminosity metallicity dependence $Z_W=-0.217\pm0.046$ mag/dex was adopted from the global fit to the distance ladder in~\citet{2022ApJ...934L...7R}. Differences among direct measurements of the metallicity intercept impact $\gamma$ in the recent literature\footnote{We distinguish here between $Z_W$ as a nuisance parameter included in the distance ladder, and $\gamma$ as the direct measurement of the metallicity effect on the Leavitt law intercept determined by spectroscopy and geometric distances alone.}~\citep[e.g.,][]{2022ApJ...939...89B, 2023ApJ...955L..13B, 2024A&A...690A.246T, 2024A&A...683A.234B, 2025arXiv250522512K, 2025arXiv250715936B} were discussed in detail and determined to be explained by specifics in methodology, without relevance for the distance network due to the limited range of $\mathrm{[O/H]}$ among SN-host Cepheids~\citep[see Fig.\,21 from][]{2022ApJ...934L...7R}. A variant using $Z_W=0$ and excluding low-metallicity anchors (LMC and SMC) was introduced to assess sensitivity to uncertainties in $Z_W$ and possible mismatches between oxygen abundances determined using optical spectra and $R_{23}$.
As described in \citet{2022ApJ...934L...7R}, statistical corrections for flux contributions due to cluster Cepheids occurring in hosts \citep{2018ApJ...861...36A} were included for HST Cepheids as a background effect. Systematic corrections for time dilation were applied to measured periods of Cepheids in hosts \citep{2019A&A...631A.165A,2022A&A...658A.148A}.

\paragraph{Assumptions.}

It is assumed that the period-luminosity relation of Cepheids does not depend on host galaxy properties other than metallicity. While the period-luminosity relation can be qualitatively explained by stellar models, it is empirically calibrated and does not rely on specific assumptions about stellar physics. The period–luminosity relation is assumed to be linear over the period range $0.4 < \log P < 2$ (with period $P$ in days), and the slope of the period-luminosity relation is considered to be independent of metallicity. There might be a mild metallicity dependence for the slope but it has not been detected to be significant at more than $1\sigma$ level \citep{2024A&A...690A.246T}.

\paragraph{Connection to the Network.}
Geometric distances of Cepheids in anchor galaxies are used to calibrate the period-luminosity relation: GEDR3 parallaxes in the Milky Way, DEBs distances in the LMC and SMC, and the maser distance to NGC 4258. Cepheids in hosts are then used to calibrate the respective secondary distance indicators. 

\subsubsection{Tip of the Red Giant Branch} \label{trgb-data}

\paragraph{Overview.}
The tip of the red giant branch, as a distance indicator, relies on a well-known astrophysical phenomenon: the onset of helium burning in the electron-degenerate helium core of an old, low-mass star~\citep{2017A&A...606A..33S}. The ignition in degenerate conditions results in a runaway process which leads to the helium flash. The core mass at helium ignition varies negligibly over a wide range of stellar population ages, and in certain passbands, metallicity effects are also small or can be standardized \citep{2017A&A...606A..33S, 2018SSRv..214..113B, 2019ApJ...880...63M}. 

\paragraph{Data Sources.}
We use TRGB distances derived from a number of sources and teams as follows: 
\begin{itemize}
\item JWST -- TRGB -- SH0ES --
The SH0ES {JWST} TRGB distances are drawn from the eight hosts of 10 SN~Ia in Table 2 of~\citep{2024ApJ...976..177L}, which are based on observations from {JWST} Cycles 1 \& 2 GO-1685 and GO-2875 (PI: A. Riess). These observations were conducted in the {JWST} {F090W} and {F150W} filters. TRGB reference magnitudes were measured using a Sobel-filter based approach, anchored to the NGC 4258 measurement from~\citet{2024ApJ...966...89A}, also using data from {JWST} Cycle 1 GO-1685.
\item JWST -- TRGB -- CCHP --
The CCHP JWST TRGB distances are measured for a sample of 10 galaxies, plus the geometric anchor NGC~4258~\citep{2025arXiv250311769H, 2025ApJ...985..203F} based on observations from JWST GO–1995~\citep{2021jwst.prop.1995F}. These measurements are made in near-infrared filters: F115W is used as the primary filter, and F356W/F444W are used to measure colors. Given the notable tilt in the TRGB in these near-infrared filters, the TRGB slope is taken into account when performing the measurements.
\item JWST -- TRGB -- TRGB-SBF Project -- The TRGB-SBF Project has measured distances to fourteen nearby elliptical galaxies with JWST/NIRCam \citep{2024ApJ...973...83A, 2025ApJ...982...26A}. These TRGB measurements are anchored to NGC 4258 using JWST observations in the  F090W and F150W filters \citep{2024ApJ...966...89A}. The measurements are performed with maximum-likelihood algorithms which incorporate results from artificial star experiments, although edge-detection measurements with a Sobel filter gives very similar results, on average.
\item HST -- TRGB -- CCHP --
The tip measurements are from~\citet{2019ApJ...882...34F}, but updated with results from~\citet{2021ApJ...915...34H} that added a measurement for NGC\,5643 and determined a direct measurement for NGC\,1404.
We note that the tip measurements from~\citet{2019ApJ...882...34F} are a mix of values derived from new observations collected for the CCHP program and those derived by~\citet{2017ApJ...836...74J} and by~\citet{2019ApJ...882...34F} a single LMC-based TRGB calibration is used. 
The HST data are from the following programs: 
    GO-13691~\citep{2014hst..prop13691F}, 
    GO-15642~\citep{2018hst..prop15642F},
    GO-9351~\citep{2002hst..prop.9351R},
    GO-10497~\citep{2005hst..prop10497R}, and
    GO-10802~\citep{2006hst..prop10802R}.
For the baseline, we elect to use a single geometric anchor, NGC\,4258, for which we take the calibration from~\citet{2021ApJ...906..125J} that assumes no color-term and adopts ${M}_{{\rm{F}}814{\rm{W}}}^{{\rm{TRGB}}}=-4.050\pm 0.028$ (stat) $\pm$ 0.048 (sys) mag.
Thus, where needed we remove the TRGB zero point used in the respective paper and then adopt the NGC\,4258 result. 
\item HST -- TRGB -- EDD -- 
The TRGB measurements from \citet{2022ApJ...932...15A} provide an independent check on the CCHP TRGB measurements, using the same underlying HST data for their Type Ia supernova hosts. The photometry is performed with distinct software, and the TRGB measurements are performed with a model-fitting approach, as opposed to the edge-detection algorithms employed by the CCHP. NGC~4258 remains the sole geometric anchor, but the TRGB absolute magnitude adopted here uses an added color calibration from~\citet{2007ApJ...661..815R} to account for metallicity effects.
\item HST -- TRGB -- Li et al -- We retrieve the distances from~\citet{2025arXiv250408921L}, who adopted a Sobel filter approach to measure the TRGB in five SNe~Ia host galaxies observed by HST, as well as the result from \citet{2022ApJ...934..185D} for NGC~7814 (which is referenced within \citet{2025arXiv250408921L} and not contained in the other references above).
\end{itemize}

\paragraph{Methodological Details.}
TRGB magnitudes are identified using either edge-detection methods (e.g., Sobel filters) (cf. \citealt{2024ApJ...963L..43A} for methodological considerations) or model-fitting to the luminosity function of RGB stars. Color-based calibrations are applied to account for metallicity effects, particularly in the near-infrared. All measurements are anchored to the geometric distance to NGC 4258, with differing treatments of slope and zero point depending on filter set and methodology.

\paragraph{Systematic Corrections.}
Some teams adopt a color-based metallicity calibration to standardize the TRGB magnitude. When a different zero point was originally assumed, it is replaced by the NGC 4258-based calibration to ensure consistency.

\paragraph{Assumptions.} The different authors adopt distinct ways to deal with the metallicity dependence of the TRGB absolute magnitude. The general consistency between distances from the various works imply that the effect of these assumptions may be minor. 

\paragraph{Connection to the Network.} The TRGB distances are geometrically anchored to NGC 4258 in both the native HST and JWST systems. These TRGB distances then provide the ability to calibrate other secondary distance indicators, such as for example the magnitudes of Type Ia supernovae or Surface Brightness Fluctuations.


\subsubsection{Mira distances}

\paragraph{Overview.}
Mira variables are fundamentally-pulsating Asymptotic Giant Branch stars that, similar to Cepheids, follow a period-luminosity Relation.  Despite not being as widely-used or as well-understood as Cepheids, Miras still have two major advantages as independent calibrators of Type Ia SNe: 1) they have larger near-infrared amplitudes than Cepheids, allowing them to be characterized and detected using only near-infrared observations and 2) they are ubiquitous since they have low-to-intermediate masses (stars with initial masses ranging from $0.8 M_\odot < M < 8 M_\odot$ are expected to undergo a Mira phase of evolution), allowing them to be detected in galaxies not reachable by Cepheids. Miras are as bright as Cepheids in the near-infrared and brighter than all of the other commonly-used stellar distances indicators such as Tip of the Red Giant Branch, J-region Asymptotic Giant Branch, and RR Lyrae. 

\paragraph{Data Sources.}
Data and distance measurements to NGC 4258, NGC 1559, and M101 using Miras originate from~\citet{2018ApJ...857...67H, 2020ApJ...889....5H, 2024ApJ...963...83H} respectively. All three galaxies were observed with HST WFC3/IR, primarily in the \emph{F160W} bandpass. NGC 4258 and NGC 1559 also have \emph{F125W} and \emph{F110W} observations used to obtain colors. 

\paragraph{Methodological Details.}
Mira periods are determined using at least 10--12 epochs of photometry spaced unevenly over a minimum baseline of one year. In addition to applying cuts to remove non-variable objects, candidate O-rich Miras are selected on the basis of their peak-to-trough amplitude ($0.4$ mag $ \lesssim A \lesssim 0.8$ mag), period of variability ($200 \lesssim  P \lesssim 400$ days), and approximate color.  These cuts are intended to limit the contamination from Carbon-rich Miras and other types of variables and Asymptotic Giant Branch stars. In each galaxy, the lower period bound is set by an empirically-determined completeness limit and magnitude of contamination from Carbon-rich Miras is modeled empirically using a mixture model based on the OGLE (Optical Gravitational Lensing Experiment) Large Magellanic Cloud observations of long-period variables. Although Miras in Milky Way star clusters and the LMC have also been used as anchors in the literature~\citep[e.g.][]{2025ApJ...990...63B}, these were not adopted here to mitigate systematic uncertainties related to ground-based to HST photometric transformations. 

\paragraph{Systematic Corrections.} None.

\paragraph{Assumptions.} 
It is assumed that Oxygen-rich Miras with periods below approximately 400 days follow a linear period-luminosity Relation. The slope and zeropoint of this relation, which are empirically calibrated, as well as the amplitudes of the Miras themselves, are not expected to evolve as a function of the host environment. See~\citet{2024ApJ...963...83H} for further discussion.

\paragraph{Connection to the Network.} The zeropoint of the Mira period-luminosity Relation is calibrated using the water megamaser distance to NGC 4258. This relation is then used to determine Mira distances to NGC 1559 (host of SN 2005df) and M101 (host of SN 2011fe) and calibrate secondary distance indicators in those hosts. 


\subsubsection{J-region AGB method}

\paragraph{Overview.} The J-region asymptotic giant branch (JAGB) method is a statistical method that uses evolved stars that are redder than the red giant branch in color-magnitude diagrams.
The stars are identified using a color and magnitude selection and the application of a statistic such as mean, median, or mode. 

\paragraph{Data Sources.}
We use JAGB distances derived from the following sources:
\begin{itemize}
\item JWST -- JAGB -- SH0ES --
We draw our data for the SH0ES JAGB host galaxy distance measurements from~\citet{2024ApJ...966...20L} and~\citep{2025arXiv250205259L}, which are based on observations from both {JWST} Cycle 1 \& 2 programs GO-1685 and GO-2785 (PI: A. Riess). These observations were obtained in the {JWST} NIRCAM \emph{F150W} and \emph{F277W} filters. 
\item JWST -- JAGB -- CCHP --
We use distances from~\citet{2025ApJ...985..203F} for the CCHP JAGB distances. These distances are based on {JWST} NIRCAM observations from program GO-1995 (PI: W. Freedman) in the \emph{F115W}, \emph{F356W}, and \emph{F444W} filters.
\end{itemize}

\paragraph{Methodological Details.} The JAGB reference magnitude can be measured using the mode, mean, median, or a model fit, as currently used in the literature. To account for systematic differences in selecting any one particular statistic,~\citep{2024ApJ...966...20L} adopted the middle, or median, of the JAGB reference magnitudes across several measurement variants and use the spread in JAGB magnitude across these variants as a systematic uncertainty. The CCHP distances use exclusively the mode.

\paragraph{Systematic Corrections.} For the SH0ES dataset, crowding corrections were applied using artificial star tests, as described in~\citet{2024ApJ...966...20L} and are measured in the `outer disk,' where the crowding bias is less than 0.05~mag in both filters used for the measurement. The values from the CCHP dataset do not apply crowding corrections but seeks to minimize the effects of crowding by also making the measurement in the outer disk. 

\paragraph{Assumptions.} The JAGB method, as  currently described in the literature, assumes that the population of stars in the J-region are homogeneous such that the mean magnitude (or alternatively the median or mode) for all stars in this region remain constant. However, this may not be the case, as noted by~\citet{2020MNRAS.495.2858R, 2021MNRAS.501..933P,2023MNRAS.522..195P, 2024ApJ...966...20L, 2025arXiv250205259L}; such effects can manifest via asymmetry of the J-region luminosity function, which causes difference between the JAGB magnitude measured using different statistics. Systematics from this asymmetry can be quantified through the systematic uncertainty adopted via measurement variants. 

\paragraph{Connection to the Network.} 
The JAGB is calibrated using the water maser distance in NGC 4258 and is used to further calibrate other secondary distance indicators in the respective hosts.

\subsection{Supernovae}
\subsubsection{Type-Ia Supernovae}

\paragraph{Overview.}
Type Ia supernovae (SNe~Ia) are thermonuclear explosions of white dwarf stars in binary systems. They are among the best-established standardizable candles in cosmology due to their consistent intrinsic luminosities after empirical corrections. This allows them to serve as precise distance indicators at cosmological scales. In the local distance ladder, SNe~Ia constitute the third rung, calibrated by nearby objects with independent distance measurements.

\paragraph{Data Sources.}
We use SNe~Ia distances derived from the following sources:
\begin{itemize}
    \item \textit{Pantheon+ (optical):} 
    The Pantheon+ sample of 1701 light curves of 1550 spectroscopically confirmed SNe~Ia compiled from 18 different surveys. The sample itself and light curves are detailed in~\cite{2022ApJ...938..113S}, is cross-calibrated in~\cite{2022ApJ...938..111B}, and host-galaxy assignments and redshifts are updated in~\cite{2022PASA...39...46C}. Pantheon+ utilize the SALT2 SN Ia model from~\cite{2021MNRAS.504.4111T} and apply bias corrections to the observed light-curve quantities following~\cite{p21,p22}. The entire analysis pipeline for Pantheon+ is automated in~\cite{Hinton:2020yik} which facilitated the systematic uncertainties described in~\cite{2022ApJ...938..110B} that are utilized for this work. 
   
    \item \textit{CSP-I\&II (optical+NIR):} The Carnegie Supernova Project (CSP; \citealt{2006PASP..118....2H}) obtained high-quality light-curves of SNe~Ia in $uBgVriYJH$ bands from two observing campaigns between 2004 and 2015. The first observing campaign (CSP-I) ran from 2004 to 2009, and the second campaign (CSP-II) ran from 2011 to 2015. Light-curves from CSP-I are published in~\cite{2017AJ....154..211K}, and spectroscopic analysis is described in~\cite{2013ApJ...773...53F}. While CSP-I followed up SNe~Ia mostly from targeted search, CSP-II followed SNe~Ia that are discovered from rolling search, primarily from La Silla Quest \citep{2013PASP..125..683B}. CSP-II is described in~\cite{2019PASP..131a4002H} and in~\cite{2019PASP..131a4001P}. Spectra of SNe~Ia are published in~\cite{2024ApJ...967...20M}, and individual light-curves will be published (N. Suntzeff et al. 2026, in preparation) soon. An extensive analysis of Hubble constant using various distance calibrators is presented in~\cite{2024ApJ...970...72U}, where light-curve fitting parameters obtained from SNooPy \citep{2011AJ....141...19B} $\texttt{max\_model}$ method are available. 
    
    \item \textit{Template-independent distances in the NIR:} A compilation of all publicly available SNe~Ia with NIR observations was presented in~\cite{2023AA...679A..95G}. These were selected for having pre-maximum data so the peak magnitude could be derived by simple Gaussian Process fitting. The sample include 19 SNe~Ia in galaxies with Cepheid-based distances from SH0ES, and 55 SNe~Ia in the Hubble flow ($z> 0.01$).
    From those, all 19 SNe~Ia in the calibrator sample have $J$-band light curves, however only 16 have light-curves with enough quality in the $H$ band. Similarly, for the Hubble-flow sample, while 52 SNe~Ia have good $J$-band light-curves, 40 have $H$-band light-curves that permit the determination of the peak-brightness. 
    For the $H_0$ determination, an intrinsic scatter of $\sigma_{int}$ of 0.125 mag was included in quadrature to the observed peak magnitude uncertainty. 
    From the baseline analysis presented in~\cite{2023AA...679A..95G}, in this paper we removed the term accounting for the SH0ES ladder in the systematic uncertainty budget (0.51 km s$^{-1}$ Mpc$^{-1}$), so the total systematic uncertainty here is 1.35 km s$^{-1}$ Mpc$^{-1}$ instead of the 1.44 km s$^{-1}$ Mpc$^{-1}$ used in~\cite{2023AA...679A..95G}.
    
    \item \textit{BayeSN (optical+NIR):} An inference of the SN~Ia ``distance,'' i.e. a normalised peak magnitude, based on optical and NIR data (the latter where available), was presented in~\citep{2023MNRAS.524..235D}. The sample included 42 SNe~Ia with Cepheid distances and 18 with TRGB distances. The SN~Ia data is compiled from the literature (individual SN references are in Tables 1 and 2). To restrict the sample to the original training set of BayeSN, we used 67 Hubble flow SNe~Ia. The uncertainties include the fit errors on each objects and the intrinsic scatter was fitted in the MCMC code used to infer $H_0$. When comparing optical only to optical + NIR, the final $H_0$ error improved by $\sim 10\%$. For the case with a $z_{\rm min} = 0.023$, the error in $H_0$ was 1.135 km\,s$^{-1}$\,Mpc$^{-1}$.
    
    \item \textit{SNe~Ia in Coma:} A sample of 13 SNe~Ia located in the Coma Cluster was recently compiled by~\citep{2025ApJ...979L...9S}, using data from multiple modern surveys including ATLAS and YSE. These SNe~Ia were selected to be hosted by passive galaxies in the Coma cluster and are located at effectively the same redshift, providing a local standard-candle ensemble. Light curves are calibrated following the Pantheon+ framework.
\end{itemize}

\paragraph{Systematic Corrections.}
Systematic uncertainties include photometric calibration, light-curve model assumptions, and corrections for selection biases. These are treated internally within each analysis and propagated into the final uncertainties on $H_0$\,.

\paragraph{Assumptions.} SNe~Ia are assumed to be standardizable candles, and the standardization is empirically-calibrated. Therefore, it is assumed that there is no evolution of SNe~Ia properties  as function of time/distance and no dependence on properties of the host such as mass or metallicity  beyond what is modeled in the standardization, see e.g., ~\cite{2023ApJ...951...22J} for a discussion. 

\paragraph{Connection to Network.} SNe~Ia are used as tracers of cosmic expansion in the distance network. 


\subsubsection{Type-II supernovae---standard candle method} \label{sec:data:snii_SC}

\paragraph{Overview.}
Type II supernova luminosities can be standardized using the Standard Candle Method (SCM; \citealt{2002ApJ...566L..63H}), an empirical technique based on correlations with photospheric velocity and color. This allows SNe II to be used as distance indicators beyond the local Universe. 

\paragraph{Data Sources.}
The Type-II supernova sample used to derive distances through the SCM is almost the same as that in~\citet{2022MNRAS.514.4620D}. This sample consists of 89 objects with $z>$0.01 from various surveys, including the Carnegie Supernova Project-I (CSP-I; \citealt{2024A&A...692A..95A}), the Lick Observatory Supernova Survey (LOSS; \citealt{2019MNRAS.490.2799D}), the Sloan Digital Sky Survey-II SN Survey (SDSS-II; \citealt{2010ApJ...708..661D}), the Supernova Legacy Survey (SNLS; \citealt{2006ApJ...645..841N}), the Dark Energy Survey Supernova Program (DES-SN; \citealt{2020MNRAS.495.4860D}), and the Subaru Hyper Suprime-Cam Survey (SSP-HSC; \citealt{2017MNRAS.472.4233D}).  Filters used include multiple optical bands (e.g., $B$, $V$, $R$, $I$, $g$, $r$, $i$).
We exclude a total of four SNe~II with Cepheid distances from the list in~\citet{2022MNRAS.514.4620D}. The distances to three supernovae (SN 1999em, SN 1999gi, and SN 2012aw) were based on Cepheid measurements using optical images. These objects were removed for consistency with our baseline, for which all the Cepheid distances were derived using NIR images. Furthermore, SN 2005ay (NGC 3938) was also removed from the sample because it used the Cepheid distance to NGC 3982---both galaxies are part of the Ursa Major groups but are located in different subgroups (South and North, respectively). We add SN 2023ixf~\citep{2025arXiv250313974Z} which exploded in M101, a galaxy for which a Cepheid distance is available~\citep{2022ApJ...934L...7R}.
Finally, the four galaxies NGC 628, NGC 5194, NGC 6946, and NGC 7793 \citep[5~SNe II used by][]{2022MNRAS.514.4620D} with TRGB measurements were re-scaled to the distances from~\citet{2022ApJ...932...15A}, which uses NGC 4258 as the TRGB zeropoint.
Therefore, to calibrate the SN~II absolute magnitude, we use a total of ten calibrators: four objects with Cepheid measurement, five SNe~II with TRGB distances, and one in NGC 4258 with geometric maser distance.

\paragraph{Methodological Details.}
We follow the methodology of~\citet{2020MNRAS.495.4860D} to derive SN velocities, magnitudes, and colors. Expansion velocities are obtained using H$_\beta$ absorption via cross-correlation between observed spectra and an SN~II template library. Light curves are modeled using hierarchical Gaussian Processes with the George Python library, allowing accurate interpolation of magnitudes and colors at fixed epochs (43 days post-explosion).


\paragraph{Systematic Corrections.}
Magnitude corrections include Milky Way extinction, K-correction, and S-correction \citep{2020MNRAS.496.3402D,2022MNRAS.514.4620D}. CMB redshifts are taken from the NASA/IPAC Extragalactic Database and corrected for peculiar velocities using the model of Carrick et al. (2015). A residual peculiar-velocity uncertainty of 250 km s$^{-1}$ is included.

\paragraph{Assumptions.}
SNeII are assumed to be standardizable candles, and the standardization is empirically-calibrated. Therefore it is assumed that there is no significant evolution of the indicator with redshift and no environment-dependent bias (e.g., host mass, metallicity).

\paragraph{Connection to the Network.}
In the distance ladder, SNe~II are used as a tracer in the Hubble flow as secondary distance indicators calibrated using primary distance indicators.

 
\subsubsection{Type-II supernovae---spectral modeling-based approach}

\paragraph{Overview.}
Type II supernovae also offer a unique opportunity for a single-step distance estimation without the use of primary distance indicators. This path makes use of the tailored Expanding Photosphere Method (tailored EPM, \citealt{2005A&A...439..671D, 2020A&A...633A..88V, 2025AA...702A..41V}), which is a significantly improved version of the classical EPM~\citep{1974ApJ...193...27K, 1994ApJ...432...42S}. The method employs spectral modeling, which allows for the precise estimation of physical parameters of the supernova at each epoch. The latest implementation of the fitting and the distance estimation is described in~\citet{2025AA...702A..41V} and~\citet{2023A&A...678A..44C}.

\paragraph{Data Sources.}
The currently published set of Type II supernovae with astrophysical modeling distances consists of 19 supernovae: ten with $z > 0.01$ which were used for the recent $H_0$ determination of~\citet{2025AA...702A..41V}, four SNe sharing two host galaxies at $z \sim 0.01$ from the SN sibling consistency check of~\citet{2023AA...672A.129C}, and five more supernovae from three additional nearby galaxies from~\citet{2023AA...672A.129C} and~\citet{2023A&A...678A..44C}. The spectra modeled in~\citet{2025AA...702A..41V} were collected by the PESSTO survey~\citep{2015A&A...579A..40S} at the New Technology Telescope of the European Southern Observatory, while~\citet{2023AA...672A.129C} used data collected across the literature, including data from~\citet{2006PASP..118....2H,2006MNRAS.372.1315S,2011ApJ...736...76R,2019ApJ...876...19S,2022ApJ...930...34T}. The observed spectra have been recalibrated to the photometry through the estimation of synthetic colours to rule out potential contaminations. The exact procedure of the analysis and the applied steps are described in~\citet{2025AA...702A..41V}.

\paragraph{Methodological Details.}
The main improvement of this technique over the traditionally used classical EPM is the use of radiative transfer modeling, which solves the outstanding issue of continuum flux dilution and the associated correction factors, which increased the systematic errors of this method in the past~\citep[for details, see][]{2005A&A...439..671D}. The radiative transfer models used for the physical parameter estimation were derived using the radiative transfer code TARDIS~\citep{2014MNRAS.440..387K, 2019A&A...621A..29V}. These models are fit to the spectra on a $\chi^2$ basis employing the emulator and fitter described in~\citet{2020A&A...633A..88V}. The estimated physical parameters are combined with the photometry to derive the absolute flux, total reddening, and the distance of the supernova following the procedure described in~\citet{2025AA...702A..41V}.

\paragraph{Systematic Corrections.}
None.

\paragraph{Assumptions.}
Since the method is based on an astrophysical model, a relatively large number of assumptions are required for modeling the physics of SNe II. The current estimation makes use of the TARDIS radiative transfer code, which assumes a spherical supernova, without time-dependent effects, but including bound-bound, bound-free transitions, and Thomson scattering as the sources of the opacity. The full setup is described by~\citet{2019A&A...621A..29V}. These assumptions hold well for a fully or almost fully ionized ejecta, and the results agree well with codes that include time-dependent effects. 
A further modeling assumption is made for the density profile of the supernova ejecta, modeled with a power-law profile, which was shown to provide adequate results in the past, but is a potential simplification.
Finally, while the method allows for the estimation of the reddening, the total-to-selective extinction ratio $R_V$ has to be assumed for the distance estimation. For the modeling, an $R_V$ of 3.1 was assumed. A more comprehensive list of assumptions and their potential impacts is available in the discussion section of~\citep{2025AA...702A..41V}.

\paragraph{Connection to the Network.}
The SN~II tailored EPM method can be used to directly determine distances to objects in the Hubble flow.  Currently, there is no overlap between the host galaxies used for the distance network and the SNe II modeling approach, however, they could serve as a calibrator in the future. An overlap between the network and the modeling of SNe II can be created by either finding new SNe~II in host galaxies, by estimating further precise standard candle distances to hosts of modelable SNe II, or by observing siblings of SNe~Ia and SNe II in the same host galaxy.  All of these options will open new avenues in the distance network and are being actively pursued.

\paragraph{Comments}
Given the assumptions used in the spectral modeling, not all Type II supernovae can be modeled in the current framework: the radiative transfer models can only fit spectra reliably in a specific epoch range, and only for sufficiently normal Type II supernovae. These can be selected in advance based on the appearance of the spectra; however, a pre-selection is needed to limit exposure to systematic effects. 

\subsection{Surface Brightness Fluctuations}

\paragraph{Overview.}
The Surface Brightness Fluctuation (SBF) technique is a statistical distance indicator that measures pixel-to-pixel variance in the unresolved stellar populations of early-type galaxies. The variance of fluctuations, normalized to the galaxy surface brightness, scales inversely with distance squared, making the galaxy appear smoother as distance increases~\citep{1988AJ.....96..807T}.
The SBF signal, which is dominated by red giant branch stars, can be reliably  measured in nearby calibration galaxies and as a tracer of universal expansion well into the Hubble flow ($\sim$100 Mpc with HST, and potentially several times farther with JWST).

\paragraph{Data Sources.}
In 2024, JWST/NIRCam observed 14 giant elliptical galaxies to simultaneously measure both their TRGB brightnesses and SBF magnitudes~\citep{2024ApJ...973...83A,2025ApJ...982...26A}, enabling a high-precision absolute calibration of the near-IR SBF distance zero point. The TRGB distances were calibrated using the geometric distance to the megamaser galaxy NGC\,4258~\citep{2024ApJ...966...89A}. This approach provides a zero point for the WFC3/IR F110W SBF distance scale that is entirely independent of the Cepheid distance scale.  
The new TRGB-based calibration was applied to SBF measurements of 61 galaxies from~\citet{2021ApJS..255...21J} and~\citet{2021ApJ...911...65B}, using updated optical $({g-z})$ color measurements, to determine the distances used here \citep{2025ApJ...987...87J}.\footnote{\url{https://github.com/jjensen-uvu/sbf_distances_2021}} The distant galaxy SBF sample spans distances up to 100 Mpc. 
In summary, the key measurements for the adopted SBF-based distances in this study come from:
\begin{itemize}
    \item JWST/NIRCam imaging of 14 elliptical galaxies used for both SBF and TRGB measurements ~\citep{2024ApJ...973...83A,2025ApJ...982...26A}. Eight of these galaxies have HST WFC3/IR SBF data, and the full set provided a distance calibration for the Fornax and Virgo clusters that connects the JWST TRGB calibration to the HST IR SBF sample of \citet{2015ApJ...808...91J}.
    \item JWST/NIRCam TRGB distance to the anchor galaxy NGC\,4258~\citep{2024ApJ...966...89A}.
    \item HST/WFC3-IR imaging of 61 early-type galaxies in the F110W filter~\citep{2021ApJS..255...21J,2021ApJ...911...65B}.
\end{itemize}

\paragraph{Methodological Details.}
SBF distances are derived from the amplitude of spatial brightness fluctuations in galaxies, quantified by fitting the Fourier spatial power spectrum with a scaled power spectrum of the image point-spread function. 
Accurate SBF relative distances are obtained using an empirical calibration adjusted for stellar population age and metallicity variations based on optical colors. 
The recent JWST/TRGB calibration of the absolute SBF magnitude, $\overline{M}$, as a function of the observed $({g - z})$ color has replaced the earlier Cepheid-based calibration, which had been used in almost all previous SBF-based distance estimates.

\paragraph{Systematic Corrections.}
The absolute fluctuation magnitude calibration depends on galaxy color, and for this study, a new calibration with updated optical colors and slope was derived \citep{2025ApJ...987...87J}. The SBF zero point is anchored to TRGB distances that were calibrated using the geometric distance to NGC\,4258. To convert distances into Hubble-flow velocities, peculiar velocity corrections were applied based on group assignments from both the Cosmic Flows 3 database~\citep{2019MNRAS.488.5438G} and the 2M++ flow model~\citep{2015MNRAS.450..317C}, following the approach of~\citet{2021ApJ...911...65B} and~\citet{2025ApJ...982...26A}. Standard corrections for foreground extinction were also applied using the prescription from~\citet{2011ApJ...737..103S}.

\paragraph{Assumptions.} 
The main assumptions behind the SBF method are that stellar population variations between galaxies are adequately accounted for by color corrections (i.e., $\overline{M}$ vs.\ $({g-z})$ color calibration), and — for the analysis presented in this work — that the TRGB stars in the spiral megamaser galaxy NGC\,4258 at a given color are representative of those in the 14 giant elliptical galaxies used to calibrate SBF.

\paragraph{Connection to the Network.}
SBF distances are calibrated by TRGB distances from JWST/NIRCam, which are anchored to the geometric distance to NGC\,4258. 
Additionally, SBF can be used to calibrate other distance indicators, thereby adding another rung to the cosmic distance ladder. The published SBF distance to one galaxy in the Coma cluster \citep{2021ApJS..255...21J} was used to provide a calibration of the Fundamental Plane (FP) method used to measure distances to thousands of elliptical galaxies~\citep{2025MNRAS.539.3627S}. SBF has also been used to provide an alternate calibration of type Ia SNe~\citep{2023ApJ...953...35G}. However, this latter path was not included in the baseline network for this study.

\paragraph{Comments and Known Issues.}
While SBF is robust for early-type galaxies, its use is limited in spiral galaxies where star formation, clumpy dust, intermediate age asymptotic giant branch and young massive stars can all affect the RGB-based SBF fluctuation signal useful for determining the distance. The accuracy of SBF distances also depends on high-quality imaging and precise knowledge of the point spread function. The SBF method is potentially sensitive to calibration systematics from $(i)$ the uncertainty in the slope of the stellar population correction arising from the limited number of calibrators and their color range, and $(ii)$ the systematic uncertainty in the zero point calibration from TRGB measured in a single anchor galaxy (NGC\,4258), and the potential differences between the TRGB brightness in that spiral galaxy and in the elliptical galaxies used to calibrate SBF.

%%%% FUNDAMENTAL PLANE %%%%
\subsection{The Fundamental Plane calibrated in Coma\label{app:data:coma}}
\paragraph{Overview.}
The Fundamental Plane (FP) is a three-dimensional scaling relation for early-type galaxies that connects their effective radius, surface brightness, and central velocity dispersion, working as a secondary distance indicator that can measure galaxy distances with $~20-25\%$ precision~\citep{1987ApJ...313...59D,1987ApJ...313...42D}. The FP relation uses two observables of elliptical galaxies that don't require any knowledge of distances, the mean surface brightness ($I_e$) and the central velocity dispersion ($\sigma_0$), to infer the physical effective radius ($R_e$), which requires distances to measure. Comparing the physical radius to the angular radius ($\theta_e$) allows for the measurement of distances. In the distance ladder, the FP acts as a secondary distance indicator: it requires calibration through absolute distance measurements to nearby galaxy clusters, after which it can trace distances to more distant early-type galaxies in the Hubble flow, making it valuable for measuring peculiar velocities and constraining cosmological parameters like $H_0$~\citep{2025MNRAS.539.3627S,2025ApJ...979L...9S} and the growth rate of cosmic structure~\citep{2020MNRAS.494.3275A, 2020MNRAS.497.1275S, 2023MNRAS.518.2436T}.

\paragraph{Data Sources.}
Building the FP relation requires combining two complementary datasets: photometric and spectroscopic observations. The photometric component draws from the DESI Legacy Imaging Surveys~\citep{2019AJ....157..168D}, which provide multi-band imaging across the DESI footprint, while the spectroscopic data come from the Early Data Release of the DESI survey~\citep{2024AJ....168...58D}. The FP analysis included 4191 early-type galaxies within the redshift range  $0 < z < 0.1$, with central velocity dispersion measurements from DESI spectra~\citep{2025MNRAS.539.3627S}. Photometric parameters, including angular effective radius and mean surface brightness, were obtained from the DESI Legacy Imaging Surveys Data Release 9.

\paragraph{Methodological Details.}
The DESI FP methodology employed a maximum likelihood approach to fit a three-dimensional Gaussian model in the parameter space defined by $\log R_e = a \log \sigma_0 + b \log I_e + c$. Central velocity dispersions were measured from DESI spectra using pPXF~\citep{2023MNRAS.526.3273C} with stellar templates from the Indo-U.S. Coudé Feed Spectral Library~\citep{2004ApJS..152..251V}. Photometric parameters were derived from DESI Legacy Survey imaging. 

\paragraph{Systematic Corrections.}
Several systematic corrections were applied to ensure robust distance measurements. For the DESI FP analysis, a photometric zero-point correction of $+0.0234$ mag was applied to northern $r$-band magnitudes to account for systematic differences between BASS and DECaLS surveys~\citep{2025MNRAS.539.3627S}. Velocity dispersion measurements were compared with SDSS pPXF values to assess potential systematic offsets, which were incorporated into the total systematic error budget rather than applied as corrections. Aperture corrections were applied using the~\citet{1995MNRAS.276.1341J} formula to convert to standard aperture sizes. A redshift cut of $z > 0.023$ was implemented to minimise peculiar velocity contamination, and peculiar velocity corrections were applied using  pvhub\footnote{\url{https://github.com/KSaid-1/pvhub}} velocity field maps. 

\paragraph{Assumptions.}
Several key assumptions are made in the application of these distance measurement methods. For the Fundamental Plane analysis, the primary assumptions include: (1) negligible evolution of the FP relation over the redshift range $0.0 < z < 0.1$, with only a modest evolution correction of $0.85z$ applied to account for stellar population aging; (2) universality of the FP relation across different galaxy environments, from field galaxies to dense cluster cores; and (3) adoption of a flat $\Lambda$CDM cosmology with fixed deceleration parameter $q_0 = -0.55$ and jerk parameter $j_0 = 1$, fitting only for the Hubble constant $H_0$ due to the limited redshift range of the data~\citep{2025MNRAS.539.3627S}. 

\paragraph{Connection to the Network.}
Within the distance ladder, the Fundamental Plane operates as a tracer that requires calibration from other distance indicators. The precise linkage is further discussed in Appendix~\ref{sec:eq_special:FP}, and is based primarily on SNe~Ia distances, with a sub-leading fixed SBF calibration also employed.  In the future it might be feasible to extend the FP approach to a full tracer, allowing for consistent modeling of expansion parameters (like $q_0$, $j_0$) and peculiar velocities within the context of the {\DN}. For now, we only rescale the {\Hcst} values of \citet{2025MNRAS.539.3627S,2025ApJ...979L...9S}, see Appendix~\ref{sec:eq_special:FP}.

%%%% TULLY-FISHER %%%%
\subsection{Tully-Fisher relation}

\paragraph{Overview.}
The Tully-Fisher (TF) relation is an empirical scaling relation connecting the rotational velocity of spiral galaxies to their intrinsic luminosity, serving as a secondary distance indicator in the cosmic distance ladder. First used as a distance indicator by~\citet{1977A&A....54..661T}, the relation enables distance measurements to spiral galaxies by comparing their observed apparent magnitude to absolute magnitudes predicted from measured HI line widths. The TF relation is a secondary distance indicator that requires calibration from primary distance indicators to establish its absolute zero-point. The current largest TF catalogue is Cosmicflows-4~\citep{2020ApJ...902..145K} containing $\sim 10\,000$ spiral galaxies with HI line widths, redshifts, and photometry.

\paragraph{Data Sources.}
The primary dataset for TF measurements is Cosmicflows-4 (CF4; \citealt{2020ApJ...902..145K}), which represents the largest TF catalogue to date with approximately 10\,000 spiral galaxies distributed across the full sky, out of which we only use around 3\,400 with complete infrared photometry (as recommended by P. Boubel 2025, priv. comm.). The CF4 sample combines data from multiple surveys, incorporating HI line widths from various sources, optical photometry from the SDSS in the $i-$band, and infrared photometry from WISE in the W1-band. 

\paragraph{Methodological Details.}
The Tully-Fisher model assumed for this dataset is typically a linear relation with a quadratic term introduced at the bright end and an intrinsic scatter that decreases linearly with rotation velocity / luminosity~\citep{2020ApJ...902..145K,2024MNRAS.531...84B}. In the {\DN} implementation, the distances of the calibrators are solved for explicitly in the same way as all other calibrators.

\paragraph{Systematic Corrections.}
The TF analysis applies corrections for galaxy inclination. Magnitude selection functions are modeled as a Gaussian drop-off beyond magnitude limits to account for high flux limit effects. Peculiar velocity corrections are applied through a simultaneous fit of velocity field parameters, including bulk flow and external dipole components. 

\paragraph{Assumptions.}
The TF relation assumes that spiral galaxies follow a universal relation between rotational velocity and luminosity with no significant evolution over the redshift range used ($z<0.1$). The method assumes that HI line widths provide accurate tracers of rotational velocity after inclination corrections, and that the intrinsic scatter in the relation is primarily astrophysical rather than systematic. The peculiar velocity model assumes that large-scale motions can be described by a linear velocity field with bulk flow and external bulk flow components. 

\paragraph{Connection to the Network.}
The TF relation serves as a tracer of the Hubble flow in the {\DN}, requiring calibration from primary distance indicators.

%%% MEGAMASERS
\subsection{Megamasers in the Hubble flow}

\paragraph{Overview.}
AGN accretion disk megamasers similar to those in NGC 4258 are also observed in more distant galaxies, where they can be used to make single-step measurements of the Hubble constant~\citep{2007IAUS..242..399B}. ~\citet{2020ApJ...891L...1P} combined and updated these distance measurement.

\paragraph{Data Sources.}
The Megamaser Cosmology Project (MCP) has carried out an extensive survey of several thousand AGN~\citep{2015IAUGA..2255730B,2020ApJ...892...18K}, identifying a number of megamaser systems that are suitable for precise ($\lesssim$10\% uncertainty) geometric distance measurements.  VLBI mapping and spectral monitoring observations conducted by the MCP have to date yielded such measurements for five megamaser-hosting galaxies: UGC 3789~\citep{2009ApJ...695..287R,2010ApJ...718..657B,2013ApJ...767..154R}, NGC 6264~\citep{2013ApJ...767..155K}, NGC 6323~\citep{2015ApJ...800...26K}, NGC 5765b~\citep{2016ApJ...817..128G}, and CGCG 074-064~\citep{2020ApJ...890..118P}.  For the analyses presented in this paper, we adopt the distance and velocity measurements from~\citet{2020ApJ...891L...1P}, with the distance measurements reprocessed as constraints on distance modulus.  These galaxies span a range of distances from $\sim$50--130\,Mpc, and the MCP observations have utilized the Robert C. Byrd Green Bank Telescope, the Karl G. Jansky Very Large Array, and the Very Long Baseline Array.

\paragraph{Methodological Details.}
For each megamaser-hosting galaxy, the orbital motion of the masers is determined using a combination of spectral monitoring and VLBI mapping, which yield constraints on the sky-plane components of the masers' positions, as well as the line-of-sight components of their velocities and accelerations.  The orbital parameters for all masers are modeled alongside global parameters describing the disk morphology, the mass and location of the central black hole, its line-of-sight velocity, and the distance to the system.  The distance and velocity are combined to constrain {\Hcst}\,.

\paragraph{Systematic Corrections.}
A variety of peculiar velocity corrections have been tested for the megamaser sample, as detailed in~\citet{2020ApJ...891L...1P}. For the analyses presented in this paper, we adopt the 2M++ peculiar velocity correction as primary; one of the variant solutions use the velocity in the CMB frame, without peculiar velocity corrections. 

\paragraph{Assumptions.}
The analysis assumes that the masers execute Keplerian orbits within a thin disk and that their orbital motion is entirely governed by the gravitational influence of an enclosed point mass. It also assumes that fundamental physical constants (e.g., $G$, $c$, $\alpha$) do not vary across space, time, or environment.

\paragraph{Connection to the Network.}
The megamasers are used directly as tracers in the Hubble flow and do not require further calibration.
